\begin{textsample}{NEG Dim 2 – 2021\_c5 – Score -3.00 – 2021\_c5/enade\_2021\_qui\_a7d57fa827.txt}  \label{ex:f2_neg_141}
CAPÍTULO 5 DISTRIBUIÇÃO DOS CONCEITOS 5.1 PANORAMA NACIONAL DA DISTRIBUIÇÃO DOS CONCEITOS O Conceito Enade é um indicador de qualidade que avalia os cursos por intermédio dos desempenhos dos estudantes a partir dos resultados do Enade.
Os valores possíveis vão de 1 ( pior situação ) a 5 ( melhor situação ).
Os cursos com apenas um ou sem nenhum concluinte participante não obtêm o Conceito Enade, ficando Sem Conceito ( SC ).
Foram avaliados em Química 358 cursos, sendo 271 ( 75,7 \% ) de Licenciatura e 87 ( 24,3 \% ) de Bacharelado. \textbf{Enfatiza} -se, mais uma vez, que a diferença entre os cursos tabulados neste capítulo e no capítulo 2 corresponde aos cursos que não foram avaliados, ou seja, àqueles, em princípio, sem alunos concluintes que tivessem sido inscritos para a prova.
Na Tabela 5.1a, são apresentadas a quantidade e a distribuição de cursos de Licenciatura de Química participantes do Enade/2021, por faixa de conceito e Grande Região.
Observando -se os dados da Tabela 5.1a, nota -se que, dos 271 cursos participantes, 104 ( 38,4 \% ) classificaram -se com conceito 2, o valor modal.
Esse foi também o conceito modal na maioria das regiões : Norte ( 58,8 \% ), Nordeste ( 60,6 \% ) e Centro-Oeste ( 42,9 \% ).
O conceito 3 foi o segundo mais frequente em nível nacional ( 33,9 \%, correspondendo a 92 cursos ), sendo o conceito modal das regiões Sudeste ( 38,6 \% ) e Sul ( 42,9 \% ), e o conceito 4, o terceiro ( 18,5 \%, correspondendo a 50 cursos ).
Houve, ainda, 11 cursos ( 4,1 \% ) que receberam conceito 5, e outros 11 cursos que receberam conceito 1.
Dos 271 cursos de Química, três ( 1,1 \% ) ficaram Sem Conceito ( SC ).
A região Norte participou com 17 cursos ou 6,3 \% do total nacional.
Nessa região, dez cursos ( 58,8 \% em termos regionais ) obtiveram conceito 2, o conceito modal para a região, como já comentado.
Os conceitos 1 e 3 foram atribuídos, respectivamente, a quatro e a três cursos ( 23,5 \% e 17,6 \%, respectivamente ).
Nenhum curso ficou Sem Conceito ( SC ) ou foi avaliado com os conceitos 4 ou 5.
A região Nordeste participou com 94 cursos ou 34,7 \% do total nacional.
Desses, 57 cursos receberam o conceito 2, o que equivale a 60,6 \% do total regional, além de ser o conceito modal na região, como já comentado.
A dois cursos ( 2,1 \% ) atribuiu -se conceito 1, e a 27 cursos ( 28,7 \% ), conceito 3.
Nenhum curso ficou Sem Conceito ( SC ) ou recebeu o conceito 5, e oito cursos ( 8,5 \% ) receberam o conceito 4, nessa região.
Dos 83 cursos participantes da região Sudeste ( 30,6 \% do total nacional ), o conceito 3 foi o conceito modal, como já anteriormente assinalado, tendo sido obtido por 32 cursos ( 38,6 \% ).
O conceito 1 foi atribuído a dois cursos ( 2,4 \% ), e o conceito 2, a 14 cursos ( 16,9 \% ).
Receberam os conceitos 4 e 5, respectivamente, 27 cursos ( 32,5 \% ) e seis cursos ( 7,2 \% ).
Nessa região, dois cursos ( 2,4 \% ) ficaram Sem Conceito ( SC ).
Os 49 cursos da região Sul corresponderam a 18,1 \% do total nacional.
Como já indicado, houve predominância do conceito 3, atribuído a 42,9 \% dos cursos da região ( conceito modal ), ou seja, atribuído a 21 dos 49 cursos participantes na região Sul.
O conceito 1 foi atribuído a um curso ( 2,0 \% ), e os conceitos 2, 4 e 5, a 11 ( 22,4 \% ), a dez ( 20,4 \% ) e a cinco cursos ( 10,2 \% ), respectivamente.
Um curso ( 2,0 \% ) ficou Sem Conceito ( SC ) na região Sul.
Já dos 28 cursos participantes na região Centro-Oeste ( 10,3 \% do total nacional ), 12 ( 42,9 \% em termos regionais ) receberam conceito 2, o conceito modal, como já fora destacado.
Dois cursos ( 7,1 \% ) obtiveram o conceito 1, nove cursos ( 32,1 \% ) obtiveram conceito 3, e outros cinco ( 17,9 \% ), conceito 4.
Nessa região, nenhum curso ficou Sem Conceito ( SC ) ou recebeu o conceito 5.
No Gráfico 5.1a, é apresentada a distribuição cumulativa do Conceito Enade dos cursos de Licenciatura segundo Grande Região.
Em um gráfico de distribuição cumulativa, em cada valor de abscissa ( eixo x ), o valor apresentado no eixo das ordenadas ( eixo y ) é o valor acumulado da variável ( Conceito Enade ).
Observando a linha vertical pontilhada, que ajuda a visualizar o percentual de conceitos acumulados até o valor 3, ou seja, conceitos 1, 2 e 3, pode -se notar que, por exemplo, a região Norte ( linha azul ) apresenta 100 \% dos cursos com conceito 3 ou abaixo deste valor.
A região Sudeste ( linha vermelha ), por outro lado, apresenta 59,3 \% dos cursos na mesma situação : com conceito 3 ou abaixo.
Destaca -se que em algumas regiões o conceito 5 pode não ter sido alcançado por nenhum curso.
É o que ocorre na região Norte ( linha azul ), na qual todos os cursos têm conceito menor ou igual a 3, e a poligonal alcança 100 \% no conceito 3, e nas regiões Nordeste ( linha verde ) e Centro- Oeste ( linha laranja ), nas quais todos os cursos têm conceito menor ou igual a 4, e a poligonal alcança 100 \% no conceito 4.
Em linhas gerais, em um gráfico de distribuição cumulativa, poligonais mais à esquerda representam Grandes Regiões com pior distribuição, e poligonais mais à direita, Grandes Regiões com melhores conceitos.
As regiões Norte ( linha azul ), Nordeste ( linha verde ) e Centro-Oeste ( linha laranja ) apresentam os cursos com conceitos mais baixos, ao passo que as regiões Sudeste ( linha vermelha ) e Sul ( linha rosa ) apresentam a distribuição com valores maiores.
Na Tabela 5.1b, são apresentadas a quantidade e a distribuição de cursos de Bacharelado de Química participantes do Enade/2021, por faixa de conceito e Grande Região.
Observando -se os dados da Tabela 5.1b, nota -se que, dos 87 cursos participantes, 27 ( 31,0 \% ) classificaram -se com conceito 3, o valor modal.
Esse foi também o conceito modal nas regiões Norte ( 100 \% ), Nordeste ( 42,9 \%, ao conceito 2 também foram atribuídos 42,9 \% dos cursos ), Sudeste ( 28,6 \% ) e Sul ( 26,1 \%, ao conceito 2 também foram atribuídos 26,1 \% dos cursos ).
O conceito 2 foi o segundo mais frequente em nível nacional ( 28,7 \%, correspondendo a 25 cursos ), sendo esse o conceito modal da região Centro-Oeste ( 42,9 \% ).
O conceito 4 foi recebido por 14 cursos ( 16,1 \% ), o conceito 1, por 12 cursos ( 13,8 \% ) e o conceito 5, por cinco cursos ( 5,7 \% ).
Houve, ainda, quatro cursos ( 4,6 \% ) que ficaram Sem Conceito ( SC ).
A região Norte participou com um curso, 1,1 \% do total nacional.
O conceito atribuído foi o 3.
A região Nordeste participou com 14 cursos ou 16,1 \% do total nacional.
Nessa região, seis cursos ( 42,9 \% em termos regionais ) obtiveram o conceito 2 e outros seis, o conceito 3, como já comentado.
Os dois cursos restantes obtiveram o conceito 4.
Dos 42 cursos participantes da região Sudeste ( 48,3 \% do total nacional ), o conceito 3 foi o conceito modal, como já anteriormente assinalado, tendo sido obtido por 12 cursos ( 28,6 \% ).
O conceito 1 foi atribuído a sete cursos ( 16,7 \% ), e o conceito 2, a dez cursos ( 23,8 \% ).
Receberam os conceitos 4 e 5, respectivamente, oito cursos ( 19,0 \% ) e quatro cursos ( 9,5 \% ).
Nessa região, um curso ( 2,4 \% ) ficou Sem Conceito ( SC ).
Os 23 cursos da região Sul corresponderam a 26,4 \% do total nacional.
Como já indicado, houve predominância dos conceitos 2 e 3, atribuído a 26,1 \% dos cursos cada, ou seja, seis cursos receberam conceito 2, e outros seis cursos, o conceito 3.
O conceito 1 foi atribuído a cinco cursos ( 21,7 \% ), o conceito 4, a três cursos ( 13,0 \% ) e o conceito 5, a um curso ( 4,3 \% ).
Dois cursos ( 8,7 \% ) ficaram Sem Conceito ( SC ).
Já dos sete cursos participantes na região Centro-Oeste ( 8,0 \% do total nacional ), três ( 42,9 \% em termos regionais ) receberam conceito 2, o conceito modal.
Dois cursos ( 28,6 \% ) obtiveram conceito 3 e um curso ( 14,3 \% ) obteve conceito 4.
Um curso ficou Sem Conceito ( SC ), e nenhum curso recebeu o conceito 1 ou o conceito 5.
No Gráfico 5.1b, é apresentada a distribuição cumulativa do Conceito Enade dos cursos de Bacharelado segundo Grande Região.
Em um gráfico de distribuição cumulativa, em cada valor de abscissa ( eixo x ), o valor apresentado no eixo das ordenadas ( eixo y ) é o valor acumulado da variável ( Conceito Enade ).
Observando a linha vertical pontilhada, que ajuda a visualizar o percentual de conceitos acumulados até o valor 3, ou seja, conceitos 1, 2 e 3, pode -se notar que, por exemplo, a região Norte ( linha azul ) apresenta 100 \% dos cursos com conceito 3.
A região Sudeste ( linha vermelha ), por outro lado, apresenta 70,7 \% dos cursos na mesma situação : com conceito 3 ou abaixo.
Destaca -se que em algumas regiões o conceito 5 pode não ter sido alcançado por nenhum curso, como ocorre na região Norte ( linha azul ), na qual o curso tem conceito igual a 3, e a poligonal alcança 100 \% no conceito 3, e nas regiões Nordeste ( linha verde ) e Centro-Oeste ( linha laranja ), nas quais todos os cursos têm conceito menor ou igual a 4, e a poligonal alcança 100 \% no conceito 4.
Em linhas gerais, em um gráfico de distribuição cumulativa, poligonais mais à esquerda representam Grandes Regiões com pior distribuição, e poligonais mais à direita, Grandes Regiões com melhores conceitos.
A região Norte ( linha azul ) apresenta os cursos com conceitos mais baixos, ao passo que as regiões Sudeste ( linha vermelha ) e Sul ( linha rosa ) apresentam as distribuições com valores maiores.
Cursos em IES situadas nas regiões Nordeste ( linha verde ) e Centro-Oeste ( linha laranja ) apresentam uma situação intermediária. 5.2 CONCEITOS POR CATEGORIA ADMINISTRATIVA, POR MODALIDADE DE ENSINO E GRANDE REGIÃO Na Tabela 5.2a, é apresentada a distribuição dos cursos de Licenciatura de Química participantes do Enade/2021, por Categoria Administrativa e Modalidade de Ensino, de acordo com os conceitos por eles alcançados, segundo as Grandes Regiões brasileiras.
Dos 271 cursos participantes, 234 ( 86,3 \% ) eram ministrados em Instituições Públicas, e 37 ( 13,7 \% ), em Instituições Privadas.
Quanto à Modalidade de Ensino, 33 cursos eram a Distância, e 238, Presenciais.
De acordo com as informações da Tabela 5.2a, em termos nacionais, entre Instituições Públicas e Privadas, observa -se que, dos 11 cursos avaliados com conceito 5, todos eram oferecidos em IES Públicas.
Dos 234 cursos participantes de IES Públicas, o conceito 2 foi o valor modal, atribuído a 94 cursos ( 40,2 \% ).
Dos demais cursos de Instituições Públicas participantes, 76 cursos ( 32,5 \% ) foram avaliados com conceito 3, 43 cursos ( 18,4 \% ), com conceito 4 e, como já comentado, 11 cursos ( 4,7 \% ), com conceito 5.
Nesta categoria, um curso ( 0,4 \% ) ficou Sem Conceito ( SC ), e nove cursos ( 3,8 \% ) receberam o conceito 1.
Na rede privada, o conceito modal foi 3, atribuído a 16 cursos, o correspondente a 43,2 \% dos 37 cursos da categoria.
Dos demais cursos participantes, dois ( 5,4 \% ) receberam conceito 1, e dez ( 27,0 \% ), conceito 2.
O conceito 4 foi atribuído a sete cursos ( 18,9 \% ).
Nesta Categoria Administrativa, dois cursos ( 5,4 \% ) ficaram Sem Conceito ( SC ), e nenhum curso recebeu o conceito 5.
Dos 238 cursos participantes da Modalidade Presencial, os conceitos 2 e 3 foram os valores modais, atribuídos a 85 cursos cada ( 35,7 \%, em cada conceito ).
Dos demais cursos nesta modalidade de ensino, nove cursos ( 3,8 \% ) receberam conceito 1, 46 cursos ( 19,3 \% ), conceito 4, e 11 cursos ( 4,6 \% ), conceito 5.
Nesta Modalidade de Ensino, dois cursos ( 0,8 \% ) ficaram Sem Conceito ( SC ).
Na Modalidade a Distância foram 33 cursos, \textbf{distribuídos} do conceito 1 ao 4 e com um curso ficando Sem Conceito ( SC ).
O conceito modal foi o 2 com 19 cursos ( 57,6 \% ), o segundo conceito mais frequente foi o 3 com sete cursos ( 21,2 \% ), em seguida veio o conceito 4 com quatro cursos ( 12,1 \% ) e, por último, o conceito 1 com dois cursos ( 6,1 \% ).
Reafirmando, nenhum curso na Modalidade de Ensino a Distância na Habilitação Licenciatura em Química foi avaliado com conceito 5.
Na análise por região, observa -se que, na região Norte, as Instituições Públicas participaram com 16 cursos ( 94,1 \% do total regional ), nenhum ficou Sem Conceito ( SC ) ou recebeu os conceitos 4 e 5.
O conceito modal para as Instituições Públicas na região foi 2, atribuído a nove cursos, correspondendo a 56,3 \% dessa categoria na região.
Quatro cursos foram avaliados com conceito 1, e três, com conceito 3.
As Instituições Privadas participaram com um curso na região Norte ( 5,9 \% do total regional ), ao qual foi atribuído o conceito 2.
Em relação à Modalidade de Ensino, 88,2 \% dos cursos oferecidos na região Norte eram presenciais, ou seja, 15 cursos.
Para esse conjunto, o conceito modal foi o 2 com nove cursos ( 60,0 \% ) e os conceitos 1 e 3 receberam três cursos ( 20,0 \% ) cada.
Por sua vez, a Modalidade a Distância ofereceu dois cursos ( 11,8 \% ) que foram alocados aos conceitos 1 e 2.
Na região Nordeste, a rede pública concentrou 90 dos 94 cursos participantes, o equivalente a 95,7 \% do total da região, proporção mais \textbf{elevada} do que a registrada nas demais regiões brasileiras.
Dentre os cursos oferecidos por IES Públicas no Nordeste, 55 cursos foram avaliados com conceito 2, o conceito modal.
Dois cursos obtiveram conceito 1, 26 receberam o conceito 3, e sete cursos receberam conceito 4.
Nessa combinação de categoria e região, nenhum curso ficou Sem Conceito ( SC ) ou com conceito 5.
As Instituições Privadas da região Nordeste participaram com quatro cursos ( 4,3 \% do total da região ).
Desses, dois obtiveram conceito 2, o valor modal, os outros dois receberam os conceitos 3 e 4.
Nenhum curso ficou Sem Conceito ( SC ) ou com os demais conceitos nessa região.
Na Educação a Distância na região Nordeste, 11 cursos foram avaliados, correspondendo a 11,7 \% do total regional.
A moda foi estabelecida no conceito 2, com nove cursos ( 81,8 \% ), e os conceitos 1 e 4 foram alocados a um curso ( 9,1 \% ) cada.
Não houve curso Sem Conceito ( SC ) ou com os conceitos 3 ou 5.
Na Educação Presencial foram 83 cursos ou 88,3 \% do total regional.
A moda também foi no conceito 2 com 48 cursos ( 57,8 \% ).
Nenhum curso ficou Sem Conceito ( SC ) ou recebeu o conceito 5.
Um curso ( 1,2 \% ) recebeu o conceito 1, 27 cursos ( 32,5 \% ) receberam o conceito 3 e sete cursos ( 8,4 \% ) receberam o conceito 4.
Na região Sudeste, a proporção de cursos da rede pública foi 75,9 \%, correspondendo a 63 dos 83 cursos participantes.
Nessa categoria e região, o conceito modal foi 4, atribuído a 25 cursos, a nenhum curso foi atribuído o conceito 1, e um curso ficou Sem Conceito ( SC ).
Os demais foram avaliados com conceito 2 ( oito cursos ), conceito 3 ( 23 cursos ) e conceito 5 ( seis cursos ).
Dos 20 cursos oferecidos por Instituições Privadas na região Sudeste, o conceito modal foi 3, atribuído a nove cursos.
Nessa região, um curso ficou Sem Conceito ( SC ), dois cursos receberam conceito 1, seis cursos receberam o conceito 2, e outros dois cursos receberam o conceito 4.
Nenhum curso recebeu conceito 5.
Também na região Sudeste, a maior parte dos cursos foi oferecida na Modalidade Presencial : 66 de 83 cursos.
Para esses, a moda foi o conceito 3, com 26 cursos.
Os demais cursos presenciais da região foram alocados nos conceitos 1 ( dois cursos ), 2 ( sete cursos ), 4 ( 24 cursos ) e 5 ( seis cursos ).
Ficou Sem Conceito ( SC ) um curso.
Na região Sudeste, 17 cursos foram oferecidos na Modalidade Educação a Distância.
O conceito 2 foi o modal concentrando sete cursos.
Um curso ficou Sem Conceito ( SC ), seis cursos foram alocados ao conceito 3, e três cursos, ao conceito 4.
Nenhum curso recebeu os conceitos 1 ou 5.
As Instituições Públicas concentraram 39 dos 49 cursos participantes da região Sul, 79,6 \% do total regional.
Desses, 17 cursos obtiveram conceito 3, o conceito modal.
Nesta combinação de Categoria Administrativa e Grande Região, um curso recebeu o conceito 1, dez cursos receberam conceito 2, seis, conceito 4, e ainda cinco, conceito 5.
Nenhum curso ficou Sem Conceito ( SC ).
As Instituições Privadas na região Sul participaram com dez cursos ( 20,4 \% dos cursos da região ), os conceitos 3 e 4 receberam quatro cursos cada, o conceito 2 recebeu um curso, e um curso ficou Sem Conceito ( SC ).
Nenhum curso recebeu os conceitos 1 ou 5.
Na região Sul, a quase totalidade dos cursos era presencial, 47 cursos em 49, ou 95,9 \% do total regional.
O conceito modal foi o 3 com 20 cursos, os conceitos 2 e 4 receberam 10 cursos cada.
O conceito 5 receberam cinco cursos.
Um curso ficou Sem Conceito ( SC ) e outro curso foi alocado ao conceito 1.
Os dois cursos da Educação a Distância foram avaliados com os conceitos 2 e 3.
Na região Centro-Oeste, 26 dos 28 cursos participantes eram de Instituições Públicas ( 92,9 \% \% em termos regionais ).
Desses, 12 cursos receberam o conceito 2.
Dos demais cursos, dois receberam conceito 1, sete receberam conceito 3 e cinco, conceito 4.
Nenhum curso ficou Sem Conceito ( SC ) ou recebeu conceito 5.
Quanto aos dois cursos oferecidos por Instituições Privadas na região Centro-Oeste ( 7,1 \% do total regional ), ambos apresentaram conceito 3.
A maioria dos cursos oferecidos na região Centro-Oeste era presencial, 27 dos 28 cursos.
O conceito 2 foi atribuído a 11 cursos, modal.
O conceito 3 foi atribuído a nove cursos, o 4, a cinco cursos e o 1, a dois cursos.
O curso a distância na região foi avaliado com conceito 2.
No Gráfico 5.2a, apresenta -se a distribuição do Conceito Enade dos cursos de Licenciatura, segundo a Categoria Administrativa da IES.
Os cursos em IES Públicas ( linha azul ) apresentam uma poligonal mais à direita do que os em IES Privadas ( linha verde ) e, consequentemente, uma distribuição de Conceito Enade com valores maiores.
Na Tabela 5.2b, é apresentada a distribuição dos cursos de Bacharelado de Química participantes do Enade/2021, por Categoria Administrativa e Modalidade de Ensino, de acordo com os conceitos por eles alcançados, segundo as Grandes Regiões brasileiras.
Dos 87 cursos participantes, 57 ( 65,5 \% ) eram ministrados em Instituições Públicas, e 30 ( 34,5 \% ), em Instituições Privadas.
Quanto à Modalidade de Ensino, dois cursos eram a Distância, e 85, Presenciais.
De acordo com as informações da Tabela 5.2b, em termos nacionais, entre Instituições Públicas e Privadas, observa -se que, dos cinco cursos avaliados com conceito 5, quatro eram oferecidos em IES Públicas, e um, em IES Privadas.
Dos 57 cursos participantes de IES Públicas, o conceito 3 foi o valor modal, atribuído a 24 cursos ( 42,1 \% ).
Dos demais cursos de Instituições Públicas participantes, 15 cursos ( 26,3 \% ) foram avaliados com conceito 2, 14 cursos ( 24,6 \% ) foram avaliados com conceito 4 e, quatro cursos ( 7,0 \% ) receberam conceito 5.
Nenhum curso ficou Sem Conceito ( SC ) ou recebeu o conceito 1.
Na rede privada, o conceito modal foi 1, atribuído a 12 cursos, o correspondente a 40,0 \% dos 30 cursos da categoria.
Dos demais cursos participantes, dez ( 33,3 \% ) receberam conceito 2, e três ( 10,0 \% ), conceito 3.
O conceito 5 foi atribuído a um curso ( 3,3 \% ).
Nesta Categoria Administrativa, quatro cursos ( 13,3 \% ) ficaram Sem Conceito ( SC ), e nenhum curso foi avaliado com o conceito 4.
Dos 85 cursos participantes da modalidade presencial, o conceito 3 foi o valor modal, atribuído a 27 cursos ( 31,8 \% ).
Dos demais cursos nesta Modalidade de Ensino, 11 cursos ( 12,9 \% ) receberam conceito 1, 24 cursos ( 28,2 \% ) receberam conceito 2, 14 cursos ( 16,5 \% ), conceito 4, e cinco cursos ( 5,9 \% ), conceito 5.
Nesta Modalidade de Ensino, quatro cursos ( 4,7 \% ) ficaram Sem Conceito ( SC ).
Aos dois cursos a Distância informados, atribuíram -se os conceitos 1 e 2.
Na análise por região, observa -se que, na região Norte apenas um curso foi avaliado.
Segundo a Categoria Administrativa foi em IES Públicas e segundo a Modalidade de Ensino foi em Educação Presencial.
Esse curso recebeu o conceito 3.
Na região Nordeste, todos os 14 cursos eram em IES Pública e Presenciais.
A distribuição foi bimodal nos conceitos 2 e 3, atribuídos a seis cursos cada.
Aos dois cursos restantes foi atribuído o conceito 4.
Na região Sudeste, a proporção de cursos da rede pública, 54,8 \%, corresponde a 23 dos 42 cursos participantes.
Nessa categoria e região, o conceito modal foi 3, atribuído a nove cursos, e nenhum curso ficou Sem Conceito ( SC ) ou recebeu conceito 1.
Os demais foram avaliados com conceito 2 ( três cursos ), conceito 4 ( oito cursos ) e conceito 5 ( três cursos ).
Dos 19 cursos oferecidos por Instituições Privadas na região Sudeste, os conceitos 1 e 2 foram atribuídos a sete cursos cada ( distribuição bimodal ).
Nessa região, um curso de IES Privadas recebeu conceito, 5 e três receberam conceito 3.
Um curso ficou Sem Conceito ( SC ), e nenhum foi avaliado com conceito 4.
Também na região Sudeste, a quase totalidade dos cursos foi oferecida na Modalidade Presencial : 41 cursos.
Para esses, a moda foi o conceito 3, com 12 cursos.
Os demais cursos Presenciais da região foram alocados nos conceitos 1 ( sete cursos ), 2 ( nove cursos ), 4 ( oito cursos ) e 5 ( quatro cursos ).
Ficou Sem Conceito ( SC ) um curso.
Na região Sudeste, o curso oferecido na modalidade Educação a Distância obteve o conceito 2.
As Instituições Privadas concentraram 10 dos 23 cursos participantes da região Sul, 43,5 \% do total regional.
Desses, cinco cursos obtiveram conceito 1, o conceito modal.
Nesta combinação de Categoria Administrativa e Grande Região, três cursos receberam conceito 2.
Nenhum curso recebeu os conceitos 3, 4 ou 5, e dois cursos ficaram Sem Conceito ( SC ).
As Instituições Públicas na região Sul participaram com 13 cursos ( 56,5 \% dos cursos da região ), seis cursos receberam conceito 3, três cursos receberam conceito 2, outros três receberam conceito 4, e um curso recebeu conceito 5.
Nenhum curso ficou Sem Conceito ( SC ) ou recebeu o conceito 1.
A quase totalidade dos cursos, 95,7 \% oferecidos na região Sul eram Presenciais.
É o mesmo que dizer que 22 dos 23 cursos da região Sul eram na Modalidade Presencial.
Os conceitos modais foram o 2 e o 3 atribuídos a seis cursos cada, os conceitos 1, 4 e 5 receberam quatro, três e um curso, respectivamente.
Finalizando, dois cursos ficaram Sem Conceito ( SC ).
O único curso na Modalidade a Distância foi avaliado com conceito 1.
Na região Centro-Oeste todos os cursos eram na Modalidade Presencial, e seis dos sete cursos eram em IES Públicas.
Desses, três cursos receberam conceito 2, dois cursos, conceito 3 e um curso, conceito 4.
O curso em IES Privadas ficou Sem Conceito ( SC ).
No Gráfico 5.2b, apresenta -se a distribuição do Conceito Enade dos cursos de Bacharelado, segundo a Categoria Administrativa da IES.
Os cursos em IES Públicas ( linha azul ) apresentam uma poligonal mais à direita do que os em IES Privadas ( linha verde ) e, consequentemente, uma distribuição de Conceitos Enade com valores maiores. 5.3 CONCEITOS POR ORGANIZAÇÃO ACADÊMICA E GRANDE REGIÃO Na Tabela 5.3a, encontra -se a distribuição dos conceitos atribuídos aos cursos de Licenciatura participantes do Enade/2021, na Área de Química, por Organização Acadêmica, segundo as Grandes Regiões brasileiras.
Dos 271 cursos de Química participantes, 170 eram oferecidos em Universidades, nove em Centros Universitários, cinco em Faculdades e 87 em CEFET/IF.
Esta distribuição corresponde a, respectivamente, 62,7 \%, 3,3 \%, 1,8 \% e 32,1 \% dos cursos.
De acordo com os dados apresentados, de todos os 11 cursos avaliados com conceito 5, dez eram vinculados a Universidades.
Os cursos vinculados a esse tipo de Organização Acadêmica tiveram o conceito 3 como conceito modal, atribuído a 58 cursos ( 34,1 \% ).
Os demais cursos vinculados a Universidades receberam os conceitos 1 ( cinco cursos ), 2 ( 57 cursos ), 4 ( 38 cursos ) e conceito 5 ( dez cursos, como já mencionado ).
Dois cursos ficaram Sem Conceito ( SC ).
Em Centros Universitários, o conceito modal foi 2, atribuído a três cursos ( 33,3 \% ).
Um curso vinculado a esse tipo de Organização Acadêmica ficou Sem Conceito ( SC ), e os outros cursos receberam os conceitos 1 ( dois cursos ), 3 ( dois cursos ) e 4 ( um curso ).
Nenhum curso recebeu conceito 5.
Dos cinco cursos mantidos por Faculdades, quatro ( 80,0 \% ) receberam conceito 3, o conceito modal, e um curso recebeu o conceito 2.
Os 87 cursos em CEFET/IF receberam conceitos.
O conceito 2 foi o modal com 43 cursos.
Os demais conceitos : 1, 3, 4 e 5 foram atribuídos a quatro, 28, 11 e um curso, respectivamente.
Considerando -se, separadamente, as regiões brasileiras, verifica -se que, na região Norte, as Universidades participaram com 12 dos 17 cursos avaliados, com conceito 2 para sete cursos e conceitos 1 e 3 para três e dois cursos, respectivamente.
Nenhum curso ficou Sem Conceito ( SC ) ou recebeu os demais conceitos nesta combinação de categoria e região.
O Centro Universitário da região Norte recebeu conceito 2.
Nenhum curso em Faculdades foi avaliado na região.
Aos quatro cursos em CEFET/IF foram atribuídos conceitos 1 ( um curso ), 2 ( dois cursos ) e 3 ( um curso ).
Na região Nordeste, as Universidades participaram com 54 dos 94 cursos da Área de Química ( Licenciatura ) oferecidos.
Nenhum desses 54 cursos ficou Sem Conceito ( SC ).
O conceito modal 2 foi atribuído a 31 cursos.
Dois cursos obtiveram conceito 1, 15 cursos obtiveram conceito 3, e seis cursos, o conceito 4.
Nenhum curso oferecido por Universidades no Nordeste recebeu conceito 5.
Os Centros Universitários contaram com dois cursos participantes na região Nordeste recebendo conceitos 3 e 4.
As Faculdades foram representadas por um curso com conceito 2.
Os CEFET/IF contaram com 37 cursos participantes na região Nordeste, 25 dos quais receberam conceito 2, o conceito modal.
Os conceitos 3 e 4 foram atribuídos a 11 e a um curso, respectivamente.
Ou seja, nenhum curso ficou Sem Conceito ( SC ) ou recebeu os conceitos 1 ou 5.
Na região Sudeste, as Universidades concentraram 57 dos 83 cursos de Química ( Licenciatura ) da região.
Dos cursos oferecidos em Universidades, nessa região, o conceito modal foi 4, atribuído a 20 cursos.
Os demais cursos receberam os conceitos 2 ( 12 cursos ), 3 ( 19 cursos ) e 5 ( cinco cursos ).
Um curso ficou Sem Conceito ( SC ), e nenhum curso obteve conceito 1.
Os Centros Universitários participaram com quatro cursos na região Sudeste, dois receberam conceito 1, um curso recebeu conceito 2, e outro curso ficou Sem Conceito ( SC ).
As Faculdades também participaram com quatro cursos nessa região, sendo que todos foram avaliados com conceito 3.
Já os CEFET/IF foram representados por 18 cursos na região Sudeste, dos quais o conceito modal 3 foi atribuído a nove cursos.
À outra metade foram atribuídos os conceitos 2 ( um curso ), 4 ( sete cursos ) e 5 ( um curso ).
Nesse tipo de organização, nenhum curso recebeu conceito 1 ou ficou Sem Conceito ( SC ).
Dos 49 cursos da região Sul, 33 eram mantidos por Universidades, para os quais o conceito modal foi 3, atribuído a 16 cursos.
Os demais cursos receberam os conceitos 2 ( três cursos ), 4 ( oito cursos ) e 5 ( cinco cursos ).
Nesse tipo de organização, um curso da região Sul ficou Sem Conceito ( SC ), e nenhum recebeu conceito 1.
Os dois cursos participantes de Centros Universitários da região Sul receberam os conceitos 2 e 3.
Não houve curso avaliado em Faculdades na região Sul.
Nos CEFET/IF, o conceito modal foi 2 atribuído a sete cursos, aos demais foram atribuídos os conceitos 1 ( a um curso ), 3 ( a quatro cursos ) e 4 ( a dois cursos ).
Na região Centro-Oeste, a metade dos 28 cursos oferecidos eram mantidos por Universidades.
Para este tipo de Organização Acadêmica, os conceitos 2 e 4 foram atribuídos a quatro cursos cada, e o conceito 3 atribuído a seis cursos ( modal ).
Nesse tipo de organização, nenhum curso recebeu os demais conceitos.
Os CEFET/IF da região Centro-Oeste contaram com a outra metade dos cursos, 14 em 28.
O conceito modal foi o 2 atribuído a oito cursos, o conceito 1, a dois cursos, o conceito 3, a três cursos, e o 4, a um curso.
Nenhum curso ficou Sem Conceito ou recebeu o conceito 5.
Nessa região não houve cursos avaliados em Centros Universitários e em Faculdades.
No Gráfico 5.3a, apresenta -se a distribuição do Conceito Enade dos cursos de Licenciatura, segundo a Organização Acadêmica da IES.
Os cursos em Universidades ( linha azul ) apresentam uma poligonal mais à direita do que as demais, denotando uma melhor distribuição de Conceitos Enade.
No outro extremo, encontram -se os cursos mantidos por Faculdades ( linha vermelha ), apresentando as respectivas poligonais mais à esquerda e, consequentemente, distribuição de Conceitos Enade com menores valores.
Os cursos mantidos por Centros Universitários ( linha verde ) e CEFET/IF ( linha laranja ) apresentam poligonais mais centrais e, consequentemente, distribuição de Conceitos Enade com valores intermediários.
Na Tabela 5.3b, encontra -se a distribuição dos conceitos atribuídos aos cursos de Bacharelado participantes do Enade/2021, na Área de Química, por Organização Acadêmica, segundo as Grandes Regiões brasileiras.
Dos 87 cursos de Química participantes, 75 eram oferecidos em Universidades, oito, em Centros Universitários, três, em Faculdades e um era oferecido em CEFET/IF.
Esta distribuição corresponde a, respectivamente, 86,2 \%, 9,2 \%, 3,4 \% e 1,1 \% dos cursos.
De acordo com os dados apresentados, todos os cinco cursos avaliados com conceito 5 eram vinculados a Universidades.
Os cursos vinculados a esse tipo de Organização Acadêmica tiveram o conceito 3 como modal, atribuído a 24 cursos ( 32,0 \% ).
Os demais cursos receberam os conceitos 1 ( oito cursos ), 2 ( 22 cursos ), 4 ( 14 cursos ) e conceito 5 ( cinco cursos, como já mencionado ).
Ficaram Sem Conceito ( SC ) dois cursos.
Em Centros Universitários, o conceito modal foi 1, atribuído a três dos oito cursos ( 37,5 \% ).
Dois cursos vinculados a esse tipo de Organização Acadêmica ficaram Sem Conceito ( SC ), e os outros cursos receberam os conceitos 2 ( dois cursos ) e 3 ( um curso ).
Aos três cursos mantidos por Faculdades, 33,3 \% cada, foram atribuídos os conceitos 1, 2 e 3.
E o curso em CEFET/IF foi avaliado com conceito 3.
Considerando -se, separadamente, as regiões brasileiras, verifica -se que, na região Norte apenas um curso foi avaliado, sendo esse em Universidade e com conceito 3.
Na região Nordeste, os 14 cursos participantes da Área de Química ( Bacharelado ) foram todos em Universidades.
A distribuição foi bimodal nos conceitos 2 e 3, avaliando seis cursos.
O conceito 4 avaliou dois cursos.
Na região Sudeste, as Universidades concentraram 32 dos 42 cursos de Química ( Bacharelado ) da região.
Dos cursos oferecidos em Universidades, nessa região, o conceito modal foi 3, atribuído a nove cursos.
Os demais cursos receberam os conceitos 1 ( quatro cursos ), 2 ( sete cursos ), 4 ( oito cursos ) e 5 ( quatro cursos ).
Nenhum curso ficou Sem Conceito ( SC ).
Os Centros Universitários participaram com seis cursos na região Sudeste, os conceitos 1 e 2 \textbf{distribuídos} a dois cursos cada, um curso recebeu conceito 3, e outro curso ficou Sem Conceito ( SC ).
As Faculdades participaram com três cursos cada um recebendo os conceitos 1, 2 e 3, respectivamente.
Os CEFET/IF participaram com um curso que recebeu conceito 3.
Dos 23 cursos da região Sul, um era mantido em Centros Universitários e recebeu conceito 1.
As Universidades participaram com 22 cursos, para os quais a distribuição foi bimodal, aos conceitos 2 e 3 foram atribuídos seis cursos cada.
O conceito 1 foi atribuído a quatro cursos, o conceito 4, a três cursos, e o conceito 5, a um curso.
Dois cursos ficaram Sem Conceito ( SC ).
Na região Centro-Oeste, os cursos avaliados na Área de Química ( Bacharelado ) eram vinculados somente a Universidades e a Centros Universitários.
Foram seis cursos em Universidades que receberam os conceitos 2 ( três cursos, modal ), 3 ( dois cursos ) e 4 ( um curso ).
E o curso em Centro Universitário ficou Sem Conceito ( SC ).
No Gráfico 5.3b, apresenta -se a distribuição do Conceito Enade dos cursos de Bacharelado, segundo a Organização Acadêmica da IES.
Nos Centros Universitários ( linha verde ), nas Faculdades ( linha vermelha ) e nos CEFET/IF ( linha laranja ), todos os cursos têm conceito 3 ou abaixo desse valor, já que a poligonal cumulativa alcança 100 \% para o conceito 3.
Os cursos em Universidades ( linha azul ) apresentam uma poligonal mais à direita do que as demais, denotando uma melhor distribuição de Conceitos Enade.
No outro extremo, com distribuição muito semelhante entre si, encontram -se os cursos mantidos por Centros Universitários ( linha verde ) e Faculdades ( linha vermelha ), apresentando as respectivas poligonais mais à esquerda e, consequentemente, distribuição de Conceitos Enade com menores valores.
A linha poligonal relativa aos CEFET/IF ( em laranja ) atinge 100 \% no conceito 3, também, mas apresenta um comportamento intermediário.

% matched lemmas: distribuir, elevado, enfatizar
\end{textsample}
